We run three further studies; Table~\ref{tab:abl} lists them. Reference values are the CNN $K{=}1$ and $K{=}4$ rows of Table~\ref{tab:skill}.

\begin{table*}[t]\centering
\caption{Sensitivity and ablation groups (mean $\pm$ std over 5 seeds). Top: rollout length $K{=}2$ (\texttt{sweep\_K}). Middle: $K{=}1$ with Gaussian input noise of std $0.03$ and $0.1$ (\texttt{abl\_noise}). Bottom: $K{=}1,4$ trained on only $4$ of the $32$ training trajectories (\texttt{sweep\_ntrain}). In each block the bold/underline ranking is only among that block's rows; the reference $K{=}1,4,8$ rows are in Tables~\ref{tab:skill} and~\ref{tab:clim}.}\label{tab:abl}
\small
\begin{tabular}{lccccccc}
\toprule
Method & rmse\_1step $\downarrow$ & rmse\_t20 $\downarrow$ & rmse\_t50 $\downarrow$ & acc\_leadtime $\uparrow$ & var\_ratio $\downarrow$ & ks $\downarrow$ & $n$ \\
\midrule
CNN K=2 & \textbf{0.023 $\pm$ 0.000} & \textbf{1.300 $\pm$ 0.034} & \textbf{3.762 $\pm$ 0.045} & \textbf{4.017 $\pm$ 0.053} & \textbf{1.001 $\pm$ 0.003} & \textbf{0.002 $\pm$ 0.001} & 5 \\
\bottomrule
\end{tabular}
\\[2pt]
\begin{tabular}{lccccccc}
\toprule
Method & rmse\_1step $\downarrow$ & rmse\_t20 $\downarrow$ & rmse\_t50 $\downarrow$ & acc\_leadtime $\uparrow$ & var\_ratio $\downarrow$ & ks $\downarrow$ & $n$ \\
\midrule
CNN K=1 noise=0.03 & \textbf{0.025 $\pm$ 0.000} & \textbf{1.564 $\pm$ 0.032} & \textbf{3.822 $\pm$ 0.024} & \textbf{3.808 $\pm$ 0.052} & \underline{0.985 $\pm$ 0.007} & \textbf{0.004 $\pm$ 0.001} & 5 \\
CNN K=1 noise=0.1 & \underline{0.086 $\pm$ 0.002} & \underline{2.770 $\pm$ 0.045} & \underline{4.230 $\pm$ 0.041} & \underline{2.467 $\pm$ 0.041} & \textbf{0.902 $\pm$ 0.013} & \underline{0.046 $\pm$ 0.009} & 5 \\
\bottomrule
\end{tabular}
\\[2pt]
\begin{tabular}{lccccccc}
\toprule
Method & rmse\_1step $\downarrow$ & rmse\_t20 $\downarrow$ & rmse\_t50 $\downarrow$ & acc\_leadtime $\uparrow$ & var\_ratio $\downarrow$ & ks $\downarrow$ & $n$ \\
\midrule
CNN K=1 ntrain=4 & \textbf{0.023 $\pm$ 0.000} & \underline{1.320 $\pm$ 0.038} & \underline{3.779 $\pm$ 0.038} & \underline{3.978 $\pm$ 0.041} & \underline{1.003 $\pm$ 0.003} & \underline{0.002 $\pm$ 0.001} & 5 \\
CNN K=4 ntrain=4 & \underline{0.024 $\pm$ 0.000} & \textbf{1.296 $\pm$ 0.027} & \textbf{3.752 $\pm$ 0.025} & \textbf{4.020 $\pm$ 0.029} & \textbf{1.002 $\pm$ 0.003} & \textbf{0.002 $\pm$ 0.001} & 5 \\
\bottomrule
\end{tabular}

\end{table*}

\textbf{Intermediate rollout length ($K{=}2$).} The $K{=}2$ model sits between $K{=}1$ and $K{=}4$ on every skill metric: RMSE at lead 2 is \rhval{sweep_k/cnn-k-2/l96/rmse_t20/mean:3} and ACC lead time \rhval{sweep_k/cnn-k-2/l96/acc_leadtime/mean:3}, with variance ratio \rhval{sweep_k/cnn-k-2/l96/var_ratio/mean:3}. The improvement from $K{=}1$ to $K{=}4$ is thus gradual, and only the largest $K$ separates clearly from noise (Fig.~\ref{fig:sweepk}). We did not test $K>8$ (cost) so we cannot say where gains saturate or reverse.

\textbf{Input noise hurts here, and is the only intervention that damps variance.} Training $K{=}1$ with input noise of std 0.03 raises RMSE at lead 2 to \rhval{abl_noise/cnn-k-1-noise-0.03/l96/rmse_t20/mean:3} (reference \rhval{main/cnn-k-1/l96/rmse_t20/mean:3}) and shortens the ACC lead time to \rhval{abl_noise/cnn-k-1-noise-0.03/l96/acc_leadtime/mean:2}; with std 0.1 the numbers are \rhval{abl_noise/cnn-k-1-noise-0.1/l96/rmse_t20/mean:3} and \rhval{abl_noise/cnn-k-1-noise-0.1/l96/acc_leadtime/mean:2}. The variance ratio falls to \rhval{abl_noise/cnn-k-1-noise-0.03/l96/var_ratio/mean:3} and \rhval{abl_noise/cnn-k-1-noise-0.1/l96/var_ratio/mean:3}, the roughness ratio to \rhval{abl_noise/cnn-k-1-noise-0.03/l96/roughness_ratio/mean:3} and \rhval{abl_noise/cnn-k-1-noise-0.1/l96/roughness_ratio/mean:3}, and the KS distance rises to \rhval{abl_noise/cnn-k-1-noise-0.1/l96/ks/mean:3} for the larger noise. All noisy runs remain stable (\texttt{stable\_frac} \rhval{abl_noise/cnn-k-1-noise-0.1/l96/stable_frac/mean:1}). Noise therefore damps small scales (the noisy model learns a smoothed map because it cannot rely on grid-scale detail of the input), which is consistent with the roughness drop; we did not isolate the mechanism further, so this is an interpretation. Noise would be expected to pay off only where the unperturbed model is unstable, which was not the case; the comparison is cross-group, so the differences are read from the table rather than from \texttt{rh compare}.

\textbf{Training-set size.} Reducing training data to $4$ trajectories (about $5000$ snapshots) leaves everything unchanged: RMSE at lead 2 is \rhval{sweep_ntrain/cnn-k-1-ntrain-4/l96/rmse_t20/mean:3} for $K{=}1$ and \rhval{sweep_ntrain/cnn-k-4-ntrain-4/l96/rmse_t20/mean:3} for $K{=}4$, against \rhval{main/cnn-k-1/l96/rmse_t20/mean:3} and \rhval{main/cnn-k-4/l96/rmse_t20/mean:3} with all $32$; both stay stable with variance ratio \rhval{sweep_ntrain/cnn-k-1-ntrain-4/l96/var_ratio/mean:3} and \rhval{sweep_ntrain/cnn-k-4-ntrain-4/l96/var_ratio/mean:3}. The ordering $K{=}4$ better than $K{=}1$ at this lead is preserved, with a similar small gap. The relevant limit is therefore not the amount of data but the information in the slow variables (unobserved fast scales) and the short training; at this data size we cannot say whether rollout training matters more at even smaller data.

\textbf{Why is one-step training stable?} We have no run that isolates a mechanism. Candidate explanations that we did not test include the very small one-step error (\rhval{main/cnn-k-1/l96/rmse_1step/mean:3} against a standard deviation of \rhval{main/cnn-k-1/l96/truth_std/mean:2}), the residual parametrisation and the dissipative attractor of Lorenz-96; they should be read as conjectures.
